% #############################################################################
% This is Chapter 10 - Conclusion
% !TEX root = ../main.tex
% #############################################################################
% PAGE BUDGET, set 2026-09-15 (see docs/thesis/REGULATIONS.md 1a-quinquies):
% 5 PAGES MAXIMUM for this chapter once all `\todo`s below are filled. A
% conclusion summarises and synthesises what earlier chapters already
% established at length - it does not re-derive or re-argue it. Summary and
% Contributions restates the RQ answers already reasoned out in
% \Cref{sec:eval_rq_answers}; Limitations and Future Work are each a
% paragraph or two of named items, not a new discussion section apiece.
% #############################################################################
\fancychapter{Conclusion}
\label{chap:conclusion}

% #############################################################################
\section{Summary and Contributions}
\label{sec:conclusion_summary}

This dissertation began from a problem that is methodological rather than technical. \ac{AI} assistance is now present in every phase of the \ac{SDLC}, adoption is close to universal and trust is not, and the evidence assembled in \Cref{chap:intro} and characterised systematically in \Cref{chap:slr} shows the value of that assistance to be conditional on the process surrounding it rather than a property of the model: functional correctness and security do not improve together, measured productivity and perceived productivity diverge, and at organisational scale output rises while delivery outcomes do not. Established delivery methodologies do not close that gap, having been designed for lifecycles in which every contributor is human, and the approaches proposed to replace them remain fragmented across isolated tools with no common position on lifecycle integration, oversight or evaluation. The research problem formalised in \Cref{chap:problem} follows directly: the absence of structured, repeatable guidance and standardised evaluation practice for integrating \ac{AI} across the \ac{SDLC} while preserving human oversight, software quality and security, and accountability.

Two artefacts were produced against that problem, and the distinction between them is load-bearing rather than presentational. The first, specified in \Cref{chap:proposal}, is a process framework: six lifecycle phases preceded by a Strategic Alignment stage, each fixing what the model may contribute, what must leave the phase, which hat is accountable and which control must be satisfied before exit (\Cref{tab:phase_mapping}); a governance set comprising Bolts as the unit of execution, the Consistency Factor, Spec-Anchored Continuity, five quality gates and three governance metrics (\Cref{sec:governance}); responsibility expressed as functions rather than positions, so that accountability survives a team of one; and three named loop-backs out of Maintenance, each re-entering a phase whose gate then applies in full. The second, presented in \Cref{chap:apex}, is Apex, a working system instantiating those constructs, built not as the contribution itself but as evidence that the constructs are implementable and as the vehicle through which the framework could be demonstrated and evaluated.

What the evaluation established is set out in \Cref{chap:demonstration,chap:evaluation} and is deliberately of two kinds. The framework was carried through a complete product over nine calendar days, growing from a pre-committed six epics to nine delivered, fifty stories and fifty-four deployed passages through all six phases, with all three loop-back routes exercised more than once each and every metric resolving to a checkable number rather than a placeholder (\Cref{sec:results_demonstration}); it was assessed analytically against five criteria fixed before the results were seen, two rated met and three partially met with the shortfall named in each case (\Cref{tab:analytical}); and it was put to five practitioners and one Agile and Scrum expert in semi-structured interviews (\Cref{sec:interviews}). The instantiation was measured separately, through a fixed nine-task script administered unmoderated to thirteen participants, returning a \ac{SUS} mean of 56.73 and a Raw \ac{NASA-TLX} grand aggregate of 21.5 out of 100. That separation is itself part of the result: a usability score for the tool is not evidence that the method is sound, and the evaluation was designed so that neither result could be quietly borrowed for the other.

The four research questions of \Cref{sec:intro_objectives} are answered in \Cref{sec:eval_rq_answers} against that evidence, and what follows summarises those answers rather than re-deriving them. \ac{DRQ}1, on mapping AI-supported activities to \ac{SDLC} phases so that the point, purpose and boundary of AI participation are unambiguous, is answered affirmatively at phase granularity: the decomposition survived instantiation with no phase merged, split or skipped, and was read back in interview by an expert who had not designed it. The qualification is that unambiguity holds at phase level and not everywhere below it, since at least one construct carried a precondition the framework never stated. \ac{DRQ}2, on the governance mechanisms, quality gates and human-in-the-loop practices that keep AI contributions traceable, validated and accountable, is answered by the set named in \Cref{sec:governance,sec:gates}, with the evidence separating those mechanisms by strength: traceability and validation are the strongest; human-in-the-loop support is met on the structural test, no phase being exitable by generation alone; and accountability is the weaker half, because the framework's own measure of AI-attributable defects is instantiable only as a proxy and a governance surface was observed reporting confidently from stale evidence.

\ac{DRQ}3, on how far those constructs can be operationalised and which resist it, is answered by the construct-by-construct mapping in \Cref{tab:apex_mapping}: of twelve constructs, ten are realised as specified, one exists only as a named proxy and one is partially realised with the resisting part identified: one metric was implemented, displayed and believed while measuring something else, caught only by comparing the implementation against its written definition, and that spec-drift detection failed because the framework assumes a specification granularity it nowhere states, so the instantiation discovered an omission in the method rather than a defect in the tool. The demonstration then showed that realised does not yet mean reliably realised: the Testing Gate was found passable by a suite that had stopped measuring the system it tested, and no gate in \Cref{tab:gates} is defined at the boundary a destructive command issued inside a phase crosses. \ac{DRQ}4, on whether the framework, through that implementation, gives practitioners usable and cognitively sustainable support, is answered split and the two halves do not contradict. On cognitive sustainability the evidence is favourable, with workload concentrated in mental demand and falling on the phases carrying the most generated content to review, which is where the framework intends to put it. On usability it is not: a grade D \ac{SUS} mean below the published corpus, error legibility the weakest questionnaire item, and nine of thirteen participants naming density or orientation as the most confusing thing about the tool. The tool asks little of a practitioner who knows where they are and does too little to get them there.

The contributions are therefore the four stated in \Cref{chap:intro}, now discharged rather than promised: the synthesis of the evidence on AI participation in the \ac{SDLC} and the systematic review supporting it; the framework itself, treating \ac{AI} as a first-class collaborator with phases, gates, function-based responsibilities, metrics and a versioned specification anchoring every artefact to the requirement it derives from; Apex, which shows those constructs implementable rather than merely describable, and which reports honestly on the two that were not; and an empirical evaluation separating what can be claimed about the method from what can be claimed about the tool. What the work does not claim is equally definite. No evidence here establishes that the framework improves delivered software quality, because no instrument in this study measures delivered quality, and the strongest results reported are perception measures over one session and one designer-operated project.

% #############################################################################
\section{Communication}
\label{sec:communication}

This section addresses the Communication activity of the \ac{DSRM}.

Dissemination follows three vehicles, each addressing a different audience and each at a different stage. The first is a scientific article reporting the \ac{SLR} of \Cref{chap:slr}, covering its protocol, the attrition from 646 retrieved articles to the 31 retained for analysis, and the five review questions it answers; the review is the component of this work that stands independently of the artefacts and is therefore the component most directly useful to other researchers. A draft exists, but submission was deprioritised in favour of completing this dissertation, and the target venue remains to be finalised.

The second is the extended abstract accompanying this dissertation, written as a self-contained technical article presenting the framework, its instantiation and the results of its evaluation to a reader who will not read the full document. It is submitted alongside this dissertation and is assessed as part of it.

The third is this dissertation itself, presented and defended before an academic examination committee. The defence is the point at which the artefacts are exposed to structured, adversarial examination by reviewers who did not design them, which no instrument in \Cref{chap:evaluation} supplies: none of its participants or interviewees was positioned to interrogate the framework's design decisions against the alternatives engaged in \Cref{sec:alternatives}. The defence is reported here as part of the \ac{DSRM} Communication activity for that reason and not merely as an academic formality.

% #############################################################################
\section{Limitations}
\label{sec:limitations}

% Keep this distinct from \Cref{sec:threats}: that section concerns the
% validity of the evaluation, this one concerns the artefact and the research
% as a whole.

The limitations below concern the artefacts and the research as a whole, kept separate from \Cref{sec:threats}, which concerns what the evaluation's design permits to be concluded from its data: a study can be internally sound and still be a study of something provisional.

\begin{itemize}[leftmargin=0.6cm, itemsep=0.5em]

\item \textbf{The framework's core constructs are an original synthesis, not an adaptation of a single validated source.} The literature grounds individual pieces, the Bolt, the Unit of Work and Mob Elaboration from \ac{AI-DLC}~\cite{AWS:2025aidlc}, the Trio's structural precedent from the Spotify model and Team Topologies~\cite{Kniberg:2012sp, Skelton:2019tt}, Spec-Anchored Continuity from spec-driven development~\cite{Thoughtworks:2025SD, Farrag:2026sd}, but no published source offers the joint blueprint assembled here, and the assembly itself has been evaluated only by the means in \Cref{chap:evaluation} (\Cref{sec:alternatives}). Two unstated premises are already established, not hypothesised: \ac{AI} Defect Escape Rate is not instantiable at the observation point its own definition presumes, since nothing says who owns production telemetry for generated logic, so only the weaker linked-post-deployment-issue proxy is realisable (\Cref{sec:apex_limitations}); and spec-drift detection presumes a per-story specification granularity the framework nowhere states, which remains open for the technical specification. The demonstration added three findings of the same kind to the two already named in \Cref{sec:conclusion_summary}'s \ac{DRQ}2 and \ac{DRQ}3 answers, the destructive-command boundary gap and the stale-instrument gate: the sharpest is that an advisory governance surface can report high confidence while wrong in either direction, with nothing distinguishing that case from being right. Whether an in-phase control can be defined without abandoning the gate model is left open.

\item \textbf{Several grounding sources are recent, vendor-reported, or not yet peer-reviewed}, identified as such at every point of use rather than laundered by citation. \ac{AI-DLC}'s productivity claims are vendor-reported and unreplicated, and its vocabulary is adopted while explicitly declining to carry its results~\cite{AWS:2025aidlc}. The principal spec-driven-development source describes the practice as emerging with unsettled scaling behaviour~\cite{Thoughtworks:2025SD}, and the strongest independent argument for specification discipline as the binding constraint on dependability is a preprint~\cite{Farrag:2026sd}. Much of \Cref{chap:intro}'s empirical case is industry telemetry and survey data rather than peer-reviewed study~\cite{DORA:2025ai, StackOverflow:2025dev, Faros:2025par, Veracode:2025gen}, and the most-cited favourable productivity figure in the field was never peer-reviewed~\cite{Peng:2023cop}. The argument rests on consistency of direction across independent sources, not the magnitude of any one figure.

\item \textbf{The vocabulary borrowed from \ac{AI} governance standards is shared language and nothing more.} \Cref{sec:metrics} maps its principles, phases, metrics and gates onto the \ac{NIST} \ac{AI} Risk Management Framework's functions~\cite{NIST:2023rmf}, with comparable vocabulary in ISO/IEC 42001:2023~\cite{ISO/IEC:2023ai}. Neither the framework nor Apex has been audited against either standard, no certification was sought, and nothing here asserts compliance with either: the claim is about where the work sits, not what it satisfies.

\item \textbf{Everything reported concerns early adoption in a single setting.} The demonstration is nine days of solo operation by the person who designed both artefacts, one project, one technology stack; the task-based study is one session of roughly an hour per participant, one seeded project, one project management tool, a convenience sample from a single professional network skewed junior; the third planned demonstration setting, an independent organisation, did not run. Nothing in that evidence base can speak to sustained use: whether gates stay exercised once the process's novelty wears off, whether the specification remains the source of truth after several amendments, whether Bolt Cycle Time and Context Traceability Rate stay informative once reported upward to someone with an interest in their value, and whether the framework survives contact with an organisation that did not choose it are all open, and none is answered here.

\end{itemize}

% #############################################################################
\section{Future Work}
\label{sec:future_work}

The directions below are derived from specific findings in this dissertation rather than assembled as a list of plausible extensions; each bullet names the finding that motivates it.

\begin{itemize}[leftmargin=0.6cm, itemsep=0.5em]

\item \textbf{Application across multiple organisations and team sizes.} The framework's accountability-by-function design is meant to survive a team of one (\Cref{sec:roles}), but the only end-to-end operation of it so far \emph{is} that limiting case: one developer holding every hat (\Cref{chap:demonstration}). The evaluation's thirteen participants ran a fixed script on a seeded project rather than their own work, and the weakest-supported \ac{UX} item was exactly the one asking whether the process fits a real team (\Cref{tab:ux_items}). Needed: application in organisations that did not design the framework, with hats actually distributed across people.

\item \textbf{Support for project management backends beyond the first-class one.} Adding a second and third adapter (\Cref{sec:apex_architecture}) showed the abstraction holds at the level of operations and fails at the level of guarantees, with three of the second tool's limitations permanent rather than unfinished. Both the demonstration and the evaluation still ran only against the first-class tool, so the parity claim is architectural, not evidenced. Needed: a demonstration against a different backend, which would also separate unfamiliarity with the tool from unfamiliarity with the board (\Cref{sec:threats}).

\item \textbf{Replacing proxy metrics with direct measurement.} \ac{AI} Defect Escape Rate's proxy only sees defects someone later linked to a story on the board; nineteen Fix-Bolts were logged across fourteen stories while the proxy itself registered one (\Cref{sec:demo_observations}). \Cref{sec:metrics} assigns this metric the job of catching a gate that has gone ceremonial, a job a same-session-blind proxy cannot do. Needed: direct measurement from production telemetry, once the framework says who owns it for generated logic specifically.

\item \textbf{Operationalising governance metrics for \ac{AI} token and compute cost.} An Agile/Scrum expert interviewed for this work (\Cref{sec:interviews}) described an organisation that centrally governs and meters \ac{AI} tool use with no role owning that governance, and asked, unprompted, whether the collaboration nets out as cost or benefit. The framework can currently answer that question for defect escape but not for cost (\Cref{sec:metrics}). Three concrete extensions, proposed the way the framework proposes its other constructs:
  \begin{itemize}[leftmargin=0.5cm, itemsep=0.3em]
  \item \emph{A fourth governance metric, \textbf{AI Token and Compute Cost}}, defined at Bolt granularity like Bolt Cycle Time (\Cref{sec:metrics}) so it composes with the existing three rather than sitting beside them: cost against Bolt Cycle Time prices flow, cost against AI Defect Escape Rate prices the defects that escaped anyway. Its observation point must be provider usage reporting at the point of invocation, named from the outset rather than discovered missing later, the way AI Defect Escape Rate's was (\Cref{sec:apex_limitations}).
  \item \emph{A named accountable hat for \ac{AI} cost governance}, plus a general mechanism for detecting an unstaffed hat. The demonstration showed the failure live: the Security Reviewer hat went unexercised through an entire Fix-Bolt cascade with nothing detecting it (\Cref{sec:demo_observations}). A gate whose evidence names the function that examined it, with a signal when that function never acted, would serve both cases at once.
  \item \emph{Risk- and phase-tiered model selection}, extending the framework's own Risk-Proportional principle, already operational in the Fast Lane/Secure Lane split (\Cref{tab:principles}), to a policy \Cref{tab:phase_mapping} does not yet fix.
  \end{itemize}

\item \textbf{Support for the Design phase specifically}, the point the evidence marks hardest to get right. Task 4 (produce and lock the design) carries the study's highest workload on every subscale but frustration and performance: mental demand 47.7 against a next-highest of 26.9, temporal demand 40.8, aggregate 30.4 (\Cref{tab:tlx_matrix}). The asymmetry is visible in the gates themselves: Testing checks generated scripts against a locked scenario that already exists; Design produces the first concrete artefact from an epic-level intent, with no earlier phase's lock to check against (\Cref{tab:gates}). Needed, not a heavier gate but a lighter one: worked examples anchoring what an adequate specification looks like for a given story type, and interim structured self-critique (usability, scenario coherence, architectural fit) before the human reviews the whole artefact at once.

\item \textbf{Redesigning the tool's interface for first-contact adoption, not sustained use.} \ac{SUS} Usability averages 59.86 against Learnability's 44.23, a full grade lower and far more dispersed (0-100 versus 12.50-87.50, \Cref{sec:results_sus}): the difficulty concentrates at first contact rather than spreading through use. Corroborated twice over: the density/orientation confusion already named in \Cref{sec:conclusion_summary}'s \ac{DRQ}4 answer above, and error legibility as the weakest individual \ac{UX} item, alongside mental demand as the largest \ac{NASA-TLX} subscale overall (\Cref{sec:results_ux}, \Cref{tab:tlx_matrix}). The Portuguese subgroup's 44.17 against the English subgroup's 60.50 points the same way, though $n=3$ cannot carry it alone. This is a direct barrier to the adoption the first future-work item above calls for, and a redesign aimed specifically at orientation and error legibility, not the tool broadly, is what the data argues for.

\end{itemize}

% #############################################################################
% Unnumbered on purpose: this is a formal declaration, not a numbered
% contribution of the Conclusion. Kept in the table of contents.
% clearpage 2026-09-13 (supervisor feedback): the declaration must be on its
% own page, not sharing a page with Future Work above it.
\clearpage
\section*{Declaration of Generative AI and AI-Assisted Technologies in the Writing Process}
\addcontentsline{toc}{section}{Declaration of Generative AI and AI-Assisted Technologies in the Writing Process}

During the preparation of this work, the author used Claude Code, with the Sonnet 5, Opus 5 and Fable 5 models, together with Google Gemini and Gemini Notebook, to improve the wording and clarity of the manuscript and to support the organisation and review of source material. After using these tools, the author reviewed and edited the text as needed and takes full responsibility for the content of the publication. The final manuscript reflects the author's own analysis, results, and conclusions.

This declaration concerns the use of these technologies in the writing process only. The use of generative AI as the object of study of this research, and within the artefacts it produces, is described in \Cref{chap:proposal} and \Cref{chap:apex}.
